\subsection{Formal power series over an integral domain}

PROPOSITION 11. --- Suppose that A is an integral domain.

\begin{enumerate}
\def\labelenumi{(\roman{enumi})}
\item
  The ring A{[}{[}I{]}{]} is again an integral domain.
\item
  If \(u, v\) are non-zero elements of \(\mathbf{A}[[\mathbf{I}]]\) , then \(\omega(uv) = \omega(u) + \omega(v)\) .
\end{enumerate}

For each \(J \subset I\) let \(\varphi_{J}\) be the homomorphism of A{[}{[}I{]}{]} into A{[}{[}J{]}{]} obtained by substituting in each element of A{[}{[}I{]}{]}, \(X_{i}\) for \(X_{i}\) when \(i \in J\) and 0 for \(X_{i}\) when \(i \in I - J\) . Let u, v be non-zero elements of A{[}{[}I{]}{]}, \(p = \omega(u)\) , \(q = \omega(v)\) ; there exists a finite subset J of I such that

\[
\varphi_ {J} (u) \neq 0, \quad \varphi_ {J} (v) \neq 0, \quad \omega (\varphi_ {J} (u)) = p, \quad \omega (\varphi_ {J} (v)) = q.
\]

Let \(a\) (resp. \(b\) ) be the homogeneous component of degree \(p\) (resp. \(q\) ) of \(\varphi_{J}(u)\) (resp. \(\varphi_{J}(v)\) ). Since \(J\) is finite, \(a\) and \(b\) are polynomials. We have \(a \neq 0\) , \(b \neq 0\) , hence \(ab \neq 0\) (IV, p.~9, Prop. 8). Hence \(\varphi_{J}(u) \varphi_{J}(v)\) is non-zero, of order \(p + q\) . It follows that \(uv \neq 0\) and \(\omega(uv) \leqslant p + q\) ; but clearly \(\omega(uv) \geqslant p + q\) .

